% GENERATED FILE. Edit canonical lessons, then run npm run generate:books.
% canonical-source-hash: fnv1a64:9508630548091fc4
% canonical-lessons: PT-C265-estava-contente, PT-C265-espero-voltar, PT-C265-revisao-sentimentos

\chapter{Estava Contente --- How It Felt, and What I Hope}
\label{ch:pt-how-it-felt-and-what-i-hope}
\hlchaptermodality{265}

\begin{chapteropening}
\textbf{By the end of this chapter:} \emph{I can say how I felt with estar in the imperfeito, and say what I hope for with espero and o meu sonho é.}

Complete PT-C265-revisao-sentimentos: tell how an experience felt and what you hope comes next.
\end{chapteropening}

\section[estava contente, estava nervoso]{estava contente — how it felt}

\label{lesson:PT-C265-estava-contente}



\begin{quote}
Say the last line of the journey you read: \textquotedblleft{}I was pleased, and so was my friend\textquotedblright{}. (\emph{Eu estava contente, e o meu amigo também.})

That sentence worked for a reason. This lesson is about why.
\end{quote}

\begin{grammarlens}[title={Grammar Lens: a feeling is scene, and it takes estar}]
You already own the feelings: \textbf{contente}, \textbf{feliz}, \textbf{triste}, \textbf{nervoso}, \textbf{preocupado}. A feeling of the moment takes \textbf{estar}, and in a story the past it wants is the imperfeito, \textbf{estava}:

\begin{quote}
\emph{Primeiro estava nervoso. Depois do exame, estava contente.} — At first I was nervous. After the exam, I was pleased.
\end{quote}

A feeling has no edge. It was going on, like the weather, so in a story it is scene. Set it against \textbf{ser}, which chapter sixteen kept for what something is by nature: \emph{ela é feliz} is a happy person; \emph{ela estava feliz} is how she was that day.

And the adjective agrees with whoever felt it. \emph{Nervoso} has four endings: \textbf{nervoso}, \textbf{nervosa}, \textbf{nervosos}, \textbf{nervosas}. \emph{Contente}, \emph{feliz} and \emph{triste} have one form for a man or a woman, and add \emph{-s} or \emph{-es} for more: \textbf{contentes}, \textbf{felizes}, \textbf{tristes}.

\begin{quote}
\emph{Estávamos contentes.} — We were pleased. \emph{A minha mãe estava preocupada.} — My mother was worried.
\end{quote}
\end{grammarlens}

\begin{grammarlens}[title={Grammar Lens: contente com}]
To say \textbf{what} you were pleased with, add \textbf{com}, \textquotedblleft{}with\textquotedblright{}:

\begin{quote}
\emph{Estava contente \textbf{com} a viagem.} — I was pleased with the journey. \emph{Estou contente \textbf{com} o meu português.} — I'm pleased with my Portuguese.
\end{quote}

\emph{Preocupado com} works the same way: \emph{estava preocupado com o exame}, worried about the exam. Portuguese says \textquotedblleft{}with\textquotedblright{} where English says \textquotedblleft{}about\textquotedblright{}.
\end{grammarlens}

\begin{cousinweb}
\textbf{contente} comes from Latin \emph{contentus}, \textquotedblleft{}satisfied\textquotedblright{}, the participle of \emph{continēre}, \textquotedblleft{}to hold together\textquotedblright{}, itself \emph{con-} plus \emph{tenēre}, \textquotedblleft{}to hold\textquotedblright{} (Lewis \& Short, under \emph{contineo}); EtymDB 2.1 traces it only as far as Old Portuguese \emph{contente}. To be content is to hold yourself in, wanting nothing further. That \emph{tenēre} is the one inside \textbf{ter} and \textbf{tinha}, and English kept both senses of the word: the \textbf{content} of a box, and a \textbf{content} person.
\end{cousinweb}

\subsection*{Your turn}
Say these aloud:

\begin{itemize}
\raggedright
  \item \emph{Primeiro estava nervoso. Depois do exame, estava contente.}
  \item the four endings — \emph{nervoso, nervosa, nervosos, nervosas}
  \item \emph{Estava contente com a viagem.}
  \item a feeling of yours, last week, with its reason — \emph{Estava … porque tinha …}
\end{itemize}

\begin{tcolorbox}[breakable,colback=teal!4,colframe=teal!35!black,title={Before you move on}]
Why is \emph{estava feliz} the usual choice in a story? (A feeling has \textbf{no edge}: it is scene, with \emph{estar}.) How many forms has \emph{contente}? (\textbf{Two} — \emph{contente}, \emph{contentes}.) Say \textquotedblleft{}I was pleased with the journey\textquotedblright{}. (\emph{\textbf{Estava contente com a viagem.}}) Which Latin verb hides in \emph{contente}? (\emph{\textbf{Tenēre}}, \textquotedblleft{}to hold\textquotedblright{}.)
\end{tcolorbox}
\section[espero voltar]{espero voltar — saying what you are reaching for}

\label{lesson:PT-C265-espero-voltar}



\begin{quote}
Say \textquotedblleft{}I hope so\textquotedblright{}. (\emph{Espero que sim.}) Then \textquotedblleft{}I was pleased\textquotedblright{}. (\emph{Estava contente.})

An experience has a past and a future. Last lesson was how it felt; this one is what you want next.
\end{quote}

\subsection*{You'll want to know: espero voltar}
\textbf{espero} + infinitive — \textquotedblleft{}I hope to\textquotedblright{}.

\begin{quote}
\textbf{Espero voltar} a Portugal. — I hope to come back to Portugal. \textbf{Espero aprender} bem português. — I hope to learn Portuguese well.
\end{quote}

You have \textbf{esperar}, which waits and hopes; \textbf{espero} is its \textquotedblleft{}I\textquotedblright{} form, and an infinitive straight after it settles the meaning: hope. Nothing goes in between. Spanish says the same, \emph{espero volver}; Italian needs a small word, \emph{spero di}; Portuguese does not.

When the person who hopes is not the one who will do it, the frame is the one you already have, \emph{espero que…}, and the verb after it bends into a form a later chapter builds. For now, keep to hopes about yourself.

\begin{grammarlens}[title={Grammar Lens: o meu sonho é}]
An ambition is bigger than a hope, and it has its own frame:

\begin{quote}
\textbf{O meu sonho é} viver em Lisboa. — My dream is to live in Lisbon. \textbf{O meu sonho é} viajar muito. — My dream is to travel a lot.
\end{quote}

\textbf{o sonho}, \textquotedblleft{}the dream\textquotedblright{}, you have. The article stays in front of \emph{meu}, as it does all through European Portuguese; in Brazil it is often dropped, \emph{meu sonho é}. After \emph{é} the infinitive follows bare: the dream \textbf{is} the doing.

Three strengths, then, for the future: \emph{gostaria de viajar}, I'd like to; \emph{espero viajar}, I hope to; \emph{o meu sonho é viajar}, it is my dream. Only the first needs a \emph{de}.
\end{grammarlens}

\begin{cousinweb}
\textbf{esperar} is inherited from Latin \emph{spērāre}, \textquotedblleft{}to hope, to expect\textquotedblright{} (EtymDB 2.1; Lewis \& Short, under \emph{spero}). English \textbf{desperate} is the same verb with \emph{dē-} in front, \textquotedblleft{}hope taken away\textquotedblright{}, and \textbf{a esperança}, \textquotedblleft{}hope\textquotedblright{}, is its noun. \textbf{o sonho} is inherited from Latin \emph{somnium}, \textquotedblleft{}a dream\textquotedblright{} (EtymDB 2.1), which is built on \emph{somnus}, \textquotedblleft{}sleep\textquotedblright{}: a dream is what sleep brings.
\end{cousinweb}

\subsection*{Your turn}
Say these aloud:

\begin{itemize}
\raggedright
  \item \emph{Espero voltar a Portugal.}
  \item \emph{Espero aprender bem português.}
  \item \emph{O meu sonho é viver em Lisboa.}
  \item the three strengths — \emph{gostaria de}, \emph{espero}, \emph{o meu sonho é} — with one verb of your own
\end{itemize}

\begin{tcolorbox}[breakable,colback=teal!4,colframe=teal!35!black,title={Before you move on}]
What goes between \emph{espero} and the infinitive? (\textbf{Nothing.}) And which of the three strengths needs a \emph{de}? (\emph{\textbf{Gostaria de.}}) Which English word is \emph{spērāre} with \emph{dē-} in front? (\textbf{Desperate.})
\end{tcolorbox}
\section[(review)]{Review — how it felt, and what comes next}

\label{lesson:PT-C265-revisao-sentimentos}



\begin{quote}
Say \textquotedblleft{}I hope to travel\textquotedblright{}. (\emph{Espero viajar.})
\end{quote}

\begin{grammarlens}[title={What you've built}]
An experience told well has two ends. The feeling looks back, with \emph{estava}; the hope looks forward, with \emph{espero} or \emph{o meu sonho é}.

\begin{quote}
\emph{A viagem era longa. Perdi o comboio, por isso estava zangado. Mas no fim estava contente com a viagem. Espero voltar.}
\end{quote}

Follow the feelings. \emph{Estava zangado} came from an event, \emph{perdi}, and \emph{por isso} joined them; \emph{estava contente} closed the story; \emph{espero voltar} opened the next one.
\end{grammarlens}

\subsection*{Your turn}
Say these aloud:

\begin{itemize}
\raggedright
  \item \emph{Perdi o comboio, por isso estava zangado.}
  \item \emph{Mas no fim estava contente com a viagem.}
  \item \emph{Espero voltar.}
  \item a journey of your own: one feeling, then a \emph{depois}, then a second feeling
  \item your dream — \emph{O meu sonho é …}
\end{itemize}

\begin{tcolorbox}[breakable,colback=teal!4,colframe=teal!35!black,title={Before you move on}]
Which verb and which past do the feelings take? (\emph{\textbf{Estar}}, in the \textbf{imperfeito}: \emph{estava zangado}, \emph{estava contente}.) Which one is an event? (\emph{\textbf{Perdi.}}) What goes between \emph{espero} and \emph{voltar}? (\textbf{Nothing.}) And why is a feeling's past called the \emph{imperfeito}? (Latin \emph{imperfectus}, \textbf{not carried through}: a feeling has no edge.)
\end{tcolorbox}
